\section{Why the adversary forces many comparisons}\label{sec:potential}

Recall from Algorithm~\ref{alg:compare} that a larger degree threshold $k$
lets the adversary retain more comparisons involving each item before
moving it into a child. Increasing $k$ lets the adversary force a comparison
count closer to $n\log n$. The following bound quantifies this dependence.

\begin{theorem}[Comparison bound]\label{thm:fixed-stage}
For every power of two $k\ge64$, the adversary of
Algorithm~\ref{alg:compare}, with the remaining parameter choices in
Table~\ref{tab:parameters}, is deterministic and legal, and every call
terminates. Every correct sorter makes at least
\begin{equation}\label{eq:exact-bound}
\frac{n\log n-3n}{1+(2\log k+17)/k}
\end{equation}
comparisons.
\end{theorem}

For this section, fix a power of two $k\ge64$.
Section~\ref{sec:runtime} chooses $k$ as a function of $n$ so that it grows
without bound while the adversary's running time remains polynomial.

The proof uses a random experiment that chooses a node and a position
for each item. At termination, all assignments admitted by the state have
the same relative order, so their probability is at most $1/n!$. We bound
this probability from below using the feasible volumes within nodes and
the probabilities of the child choices made during refinement. Relating
those child choices to the volume gained during refinement will leave
only a small additional loss per comparison.

\subsection{Volume changes}\label{sec:guarantees}

We first measure the freedom left within the individual nodes. At a node
$v$, the feasible assignments form the region
\begin{equation}\label{eq:node-region}
K_v=\{X:a_x<X_x<b_x\text{ for every free }x,
       \quad X_x<X_y\text{ for every in-node arc }x\to y\}.
\end{equation}
Write $V_v=\vol(K_v)$ and take the product over the adversary's tree:
\begin{equation}\label{eq:Z}
Z=\prod_v V_v.
\end{equation}
Each item has a position variable at exactly one node, so the product
region has dimension $n$. A node with no free items contributes volume
one. Initially $Z=1$; whenever every node has positive volume,
$0<Z\le1$.

Choose the estimation error to be
\begin{equation}\label{eq:estimation-error}
\eta=\frac1{100k}.
\end{equation}
By Lemma~\ref{lem:estimation}, each estimated answer retains at least
$1/2-\eta$ of its node's volume. Its decrease in $\log Z$ is therefore
at most
\begin{equation}\label{eq:comparison-loss}
a_\eta=-\log(1/2-\eta)\le1+\frac1{2k}.
\end{equation}
For the inequality, use
$2^{-1/(2k)}\le1-1/(4k)\le1-1/(50k)=1-2\eta$.
An answer already determined by the stored constraints leaves $Z$ unchanged.

Refinement fixes an item's position in its parent and gives it a fresh
unit coordinate in a child. Write
\begin{equation}\label{eq:refinement-gain}
g=\log\frac{Z_{\rm after}}{Z_{\rm before}}
\end{equation}
for the resulting change in $\log Z$. We call $g$ the refinement's
\emph{volume gain}; it can be negative.

The algorithm must preserve positive node volumes and the bounds needed
by the estimator. It must also return answers consistent with the earlier
comparisons. The following proposition gives these guarantees.

\begin{proposition}[Consistency and termination]\label{prop:legality}
With the parameter choices of Theorem~\ref{thm:fixed-stage}, every reached
state has positive node volumes and represents only orders consistent with
the transcript. Every comparison call terminates and performs at most $n$
refinements. Every in-node DAG has degree at most $k$ and directed paths
of fewer than $2H$ vertices throughout the interaction.
\end{proposition}

Section~\ref{sec:repairs} proves Proposition~\ref{prop:legality}.

\subsection{The cost of choosing a child}\label{sec:movement-volume}\label{sec:accounting}

The product $Z$ measures freedom within individual nodes. For example,
two items stored alone in separate children can each retain the full
interval $(0,1)$, giving $Z=1$, even though the child anchors fix their order.
To account for the order determined by the tree, we use a random experiment
on a fixed state.

Start at the root of the adversary's tree. At a visited node, stop with
probability $1/2$. Otherwise, choose an anchor by a walk in the implicit
anchor tree: at each node of that tree, select its position with probability
$1/k$; otherwise, move to the left or right child with equal probability.
If the selected anchor has an existing child in the adversary's tree,
continue the experiment in that child. Otherwise, stop at the current
node. Whenever the experiment stops in the adversary's tree, choose a
uniform coordinate in $(0,1)$ there.

The walk selecting an anchor terminates almost surely. A particular
anchor at depth $m$ requires $m$ specified left or right choices, so an
existing child at that anchor is entered with probability
\begin{equation}\label{eq:child-weight}
w(m)=\frac1{2k}\left(\frac{1-1/k}{2}\right)^m
     =\frac{2^{-\gamma m}}{2k},
\qquad \gamma=1-\log(1-1/k).
\end{equation}
The probability of stopping at each visited node of the adversary's tree
is at least $1/2$. Since the state has a finite tree, the experiment
produces a node and a coordinate almost surely.

Sample one such position independently for each of the $n$ labels.
The tree orders these positions by the same rule as in
Section~\ref{sec:overview}: compare a coordinate with the child anchors,
and order positions inside one child recursively. Ties have probability
zero, and all labels use the same distribution. Permuting the labels
therefore shows that each relative order has probability $1/n!$.

For the counting argument, measure a child choice by its contribution to
the negative logarithm of the probability:
\begin{equation}\label{eq:child-cost}
\ell(m)=-\log w(m)=\gamma m+\log k+1.
\end{equation}
The next lemma expresses the final restriction to one order in terms of
these costs.

\begin{lemma}[Bound at termination]\label{lem:packing}
Suppose the transcript determines a unique total order. If $m_r$ is the
anchor depth used by refinement $r$, then
\begin{equation}\label{eq:terminal}
\log(n!)\le n-\log Z+\sum_r\ell(m_r).
\end{equation}
\end{lemma}
\begin{proof}
Consider the event that every label stops at the node where its record is
stored, and that the sampled coordinates satisfy all the in-node
constraints. Conditional on the chosen stopping nodes, the coordinates
are independent and uniform, so the probability of satisfying the
constraints is $Z$.

The stopping probabilities contribute at least $2^{-n}$. Since records
only move downward, their required paths, counted separately for each
label, contain exactly the edge traversals made by refinements. The child
choices therefore contribute $\prod_r w(m_r)$. The event has probability
at least
\[
Z\,2^{-n}\prod_r w(m_r)=Z\,2^{-n-\sum_r\ell(m_r)}.
\]
By Proposition~\ref{prop:legality}, the event produces an order consistent
with the transcript. The unique such order has probability $1/n!$ in the
random experiment. Comparing these probabilities and taking logarithms
proves \eqref{eq:terminal}.
\end{proof}

\subsection{The guarantees needed from refinement}\label{sec:refinement-cost}

We now compare each child-choice cost with the corresponding volume gain.
The required bounds distinguish the three reasons for moving an item in
Algorithm~\ref{alg:compare}. Choose the path threshold
\begin{equation}\label{eq:path-threshold}
H=2^k.
\end{equation}
Outside a path batch, the algorithm refines a queried item because its
vertex has degree $k$ or its coordinate is narrow. When both hold, count
it as a degree refinement. Set
\begin{equation}\label{eq:charge-coefficient}
\lambda=1+\frac{\log k+10}{k}.
\end{equation}
The following lemma controls the cost of each kind of refinement using
its volume gain. Degree refinements receive an additional allowance.

\begin{lemma}[Cost of refinement]\label{lem:refinement-cost}
A degree refinement with anchor depth $m$ and volume gain $g$ satisfies
\[
\ell(m)\le\lambda g+\log k+6.
\]
A narrow refinement satisfies $\ell(m)\le\lambda g$.
For a complete path batch with depths $m_i$ and gains $g_i$,
\[
\sum_{i=1}^H\ell(m_i)\le\lambda\sum_{i=1}^H g_i.
\]
\end{lemma}

Section~\ref{sec:refinement-guarantee} proves this lemma. Narrow refinements
supply a large individual gain, and complete path batches supply a large
average gain. These gains cover their child-choice costs. Degree
refinements may need the extra allowance, which we charge to earlier
comparisons in the proof below.

\subsection{The comparison bound}\label{sec:comparison-cost}

\begin{proof}[Proof of Theorem~\ref{thm:fixed-stage}]
Proposition~\ref{prop:legality} gives consistency and termination of every
call. The procedures and their tie rules are deterministic.
Suppose the sorter finishes after $q$ comparisons, and write $j=\log k$.
Correctness requires its terminating transcript to determine a unique
total order, so Lemma~\ref{lem:packing} applies.

There are at most $q/k$ degree refinements: each deletes $k$ in-node arcs,
each comparison inserts at most one arc, and an inserted arc can be deleted
only once. Apply Lemma~\ref{lem:refinement-cost} to these refinements, the
remaining individual refinements, and the complete path batches. This gives
\begin{equation}\label{eq:refinement-total}
\sum_r\ell(m_r)\le\lambda\sum_r g_r+\frac{j+6}{k}\,q.
\end{equation}
Each comparison decreases $\log Z$ by at most $a_\eta$, while a refinement
changes it by $g_r$. Since initially $Z=1$,
\[
\sum_r g_r\le\log Z+a_\eta q.
\]
Substitute this into \eqref{eq:refinement-total} and then into the terminal
bound \eqref{eq:terminal}:
\begin{equation}\label{eq:comparison-total}
\log(n!)\le n+(\lambda-1)\log Z
              +\left(\lambda a_\eta+\frac{j+6}{k}\right)q.
\end{equation}
The term $(\lambda-1)\log Z$ is nonpositive. Moreover,
\begin{align}
\lambda a_\eta+\frac{j+6}{k}
&\le1+\frac{2j+33/2}{k}+\frac{j+10}{2k^2}\nonumber\\
&\le1+\frac{2j+17}{k},\label{eq:comparison-charge}
\end{align}
where the last inequality uses $j+10\le k$.
Finally,
\[
\log(n!)\ge\int_1^n\log x\,dx
=n\log n-\frac{n}{\ln2}+\frac1{\ln2}.
\]
Combining this estimate with \eqref{eq:comparison-total},
\eqref{eq:comparison-charge}, and $1+1/\ln2<3$ proves
\eqref{eq:exact-bound}.
\end{proof}

The next section proves the refinement and consistency guarantees used above.
